\section{The certificate: Sylvan's computation, typechecked}\label{sec:prooflog}

This section is the guide to how a goal's certificate is made and
checked. It assumes nothing from the rest of the report: every term is
defined here, with a small example, and every step can be reproduced from
the commands at the end (\cref{sec:prooflog:try}). The plan it follows is
\file{doc/certificate.md}.

\subsection{Why a certificate}\label{sec:prooflog:why}

\paragraph{Two ways to prove a fact.} Take the fact $2 + 2 = 4$ about the
natural numbers of Rocq. One way is to write a function that decides it,
for instance \code{Nat.eqb (2 + 2) 4}, which evaluates to \code{true}, and
let Rocq evaluate it: the proof is \code{eq_refl}, and Rocq's kernel accepts
it only after it has run the function itself. That is a \emph{proof by
computation} (or by reflection~\cite{boutin97}): the evidence is a
computation that the kernel repeats. The other way is to give the steps:
apply the lemma that $2 + 2 = 2 + (1 + 1)$, then the one that \dots, each
step a lemma applied to arguments, and the kernel only checks that each
application has the right type. That is a \emph{proof by certificate}: the
evidence is written down, and checking it is reading it.

\paragraph{Why it matters here.} The plugin's goals are about systems with
up to millions of states (three-process CADP: 7.9 million). The plugin
\emph{decides} a goal with a symbolic engine (\cref{sec:symbolic}), a
program in OCaml and C that is fast but not trusted. The certificate is
that program's own computation: the decision diagram library, Sylvan,
records each result it computes, the \emph{proof log}, and the plugin
writes each record as a lemma applied to its nodes. Rocq recomputes
nothing (CLAUDE.md, ``Hard rules''); \file{scripts/check_no_computation.py},
first in CI, fails on any use of \code{vm_compute}, a VM cast or the like.

\subsection{What the engine computes, from scratch}\label{sec:prooflog:background}

\paragraph{States as vectors.} The engine turns every state of a system
into a vector of natural numbers of fixed length (\cref{sec:containers}).
A pair of states, one from each system, is the two vectors put end to
end.

\paragraph{List decision diagrams.} A set of such vectors is stored as a
\emph{list decision diagram} (LDD)~\cite{blom-vdpol08}: a node holds a
value $v$, a child \emph{down} (the rest of the vectors that have $v$ at
this position) and a child \emph{right} (the vectors that have a larger
value here). Two leaves end the diagram: \code{D0}, the empty set, and
\code{D1}, the set holding only the empty vector. In Rocq
(\loc{theories/DdTree.v}{22}):
\begin{verbatim}
Inductive dd := D0 | D1 | Nd (v : N) (d r : dd).
\end{verbatim}
and \code{In_t x w} (\loc{theories/DdTree.v}{24}) says that the vector
\code{w} is in the set \code{x}: in \code{Nd v d r}, either \code{w} starts
with \code{v} and its rest is in \code{d}, or \code{w} is in \code{r}.
\emph{Example:} the set of vectors of length 1 $\{(1), (3)\}$ is
\code{Nd 1 D1 (Nd 3 D1 D0)}: a node with value 1 whose rest is the empty
vector, and to its right a node with value 3. In a node, the values along
a right chain always increase. Sylvan~\cite{vandijk-vdpol15} is the C
library that stores these nodes, each once (two equal nodes are the same
node, named by its index in Sylvan's table).

\paragraph{Operations on diagrams.} Sylvan computes with sets through a
few recursive operations: \emph{union}, \emph{minus} (difference),
\emph{intersect}, and the \emph{image} and \emph{pre-image} through a
relation (\code{lddmc_relprod}: the states one step after a set;
\code{lddmc_relprev}: the states one step before it). Each works node by
node, on the first values of its two arguments. Union, for instance
(\loc{vendor/sylvan/src/sylvan_ldd.c}{690}): if $a$ is empty, the result is
$b$ (and the other way round); if $a = b$, it is $a$; otherwise compare the
first values $v_a$ and $v_b$:
\begin{itemize}
\item $v_a < v_b$: the result is a node $v_a$, with $a$'s down, and to its
  right the union of $a$'s right with $b$;
\item $v_a > v_b$: the same with the roles swapped;
\item $v_a = v_b$: a node $v_a$ whose down is the union of the two downs
  and whose right is the union of the two rights (if the downs' union is
  empty, just the rights' union).
\end{itemize}
Sylvan remembers each result it computes in a \emph{cache}, keyed by the
operation and its arguments, so that a sub-computation met twice is done
once. (Before using the cache, its union swaps the arguments so that the
first has the larger index.)

\paragraph{The game.} To decide whether one system simulates another, the
engine solves a \emph{game} on pairs of states as a \emph{greatest
fixpoint} (\code{gfp}, \file{lib/learn/learn.ml}): start with $z$ = every
pair the engine explored (called \code{uni}); compute $\mathit{bad}(z)$,
the pairs of $z$ where the first system has a move the second cannot
answer (by silent steps, the same label, silent steps) into a pair of $z$;
remove them, $z := z \setminus \mathit{bad}(z)$; repeat until $z$ no longer
changes. The last iteration is the one where $z \setminus \mathit{bad}(z) =
z$: no pair of $z$ is bad, which is exactly what a simulation is. So
\emph{the last iteration is the proof}, if it can be written down. The
game takes 1 to 3 iterations on most goals, and up to 23 on
three-process CADP, of 6 to 97\,ms each.

The certificate does not log the game's iterations as they run: the
game's check is \emph{backward} (through the pairs from which an answer
exists, within the explored pairs \code{uni}), and a proof of it would
need \code{uni}, each system's exploration closed, and more. It logs the
same fact read \emph{forward}, as \emph{images}, which is the definition
of a simulation on $z$. For each relation $r$ of the \emph{mover} (the
system whose moves must be answered), with label $l$: every pair one
move after a pair of $z$ (the \emph{image} of $z$ by $r$) is a pair the
\emph{answerer} can answer. Those pairs, $\mathit{ok}_l$, come from $z$ by
the answerer's relations taken backwards (their \emph{inverses}: each
tuple's old and new values swapped): the pairs from which, after silent
steps, the label $l$ and silent steps, the answerer reaches $z$. The
inclusion is a difference that comes out empty. \emph{Example:}
$z = \{(p, s)\}$, the mover steps $p \xrightarrow{a} p'$ and the answerer
$s \xrightarrow{a} s'$, and $(p', s') \in z$. The image of $z$ by the
mover's $a$-relation is $\{(p', s)\}$; the image of $z$ by the answerer's
$a$-relation backwards is $\{(p', s)\}$ ($s$ steps to $s'$); and
$\{(p', s)\} \setminus \{(p', s)\}$ is empty. This check needs no
universe, no closure of the exploration, no intersection and no preimage
within a set. For a bisimilarity, each side in turn is the mover.

\subsection{The proof log}\label{sec:prooflog:log}

\paragraph{What is recorded.} While the log is open, each of the five
operations above writes one \emph{record} for each result it computes, at
the point where it offers the result to its cache
(\loc{vendor/sylvan/src/sylvan_ldd.c}{752} for union; 801, 925, 1167, 1738
for the others): which operation, which case it took (for union: whether
$v_a$ was below, equal to or above $v_b$; for the images, the kind of
position), its arguments and its result, all as node indices; 48 bytes.
Each worker thread of Sylvan writes its own buffer, so recording takes no
lock. By default the log stays in memory: each worker keeps its full
buffers in a list of its own, and the plugin reads the nodes the records
name from Sylvan's table itself. When the log is kept in a file instead,
closing it (\code{lddmc_log_close}) also writes every node reachable from
a recorded index (its value, down and right; 32 bytes), so the file can
be read without Sylvan. This is in the plugin's fork of Sylvan, on the
local branch \code{proof-log} of \file{vendor/sylvan} (three commits, on
10 and 11 November 2026; not pushed).

\paragraph{When.} The last iteration cannot be known in advance, so it is
repeated: once the fixpoint has finished, the engine clears Sylvan's cache
(else the repeated run would find every result there and record nothing),
turns Sylvan's garbage collection off (a collection would free nodes and
reuse their indices, which name the nodes in the log), grows Sylvan's
tables to their maximum first (Sylvan grows them only when it collects),
and runs, with the log open, the check as images
(\loc{lib/learn/learn.ml}{1459}). Logging every
iteration instead was rejected: the cache survives between iterations, so
every fact of every iteration would have to be kept, and CADP's 17 to 23
iterations would write many times what the proof needs.

\paragraph{Where, and whether.} Two options of the plugin
(\loc{src/user_tactics.ml}{186}) decide it:
\begin{itemize}
\item \code{Set MeBi Log Mode "off".} (the default) or \code{"keep"}:
  the log stays in memory while the tactic proves the lemma and is freed
  once the proof is built (\cref{sec:qed}), or it is written to a file and
  kept, for \code{MeBi Check Log}.
\item \code{Set MeBi Log Directory "<dir>".}: where the directory of
  the plugin's files goes; by default, next to the file being checked.
\end{itemize}
In memory, each of Sylvan's workers keeps its full buffers of records
in a list; after the log is closed, the plugin takes the records a
buffer at a time (\code{lddmc_log_take}), each buffer freed once taken,
and reads the nodes from Sylvan's own table, which collection leaves
alone until the facts are written. Nothing is written to disk. For a
file \code{F.v}, the plugin's files go in the directory \file{.F.v.mebi}
next to it (under dune, next to the copy dune checks, in
\file{_build/default}). A kept log of a lemma \code{L} is the one file
\file{.F.v.mebi/L.log}: the word \code{MEBILOG1}, the number of records,
the records, the number of nodes, the nodes. (A tactic is not told the
name of the lemma it proves; the plugin takes it from Rocq's internal
proof state, \loc{src/user_tactics.ml}{224}, and writes \code{goal} when
there is none.) \file{lib/ldd/proof_log.ml} reads such a file back.

\paragraph{Measured} (10 November 2026):

\begin{center}
\small
\begin{tabular}{@{}lrrrr@{}}
\toprule
 & records & nodes & logged check & memory \\
\midrule
two-process CADP & 10,533 & 6,234 & 0.17\,s & 0.51\,GB \\
three-process CADP & 488,067 & 233,068 & 0.82\,s & 1.33\,GB \\
\bottomrule
\end{tabular}
\end{center}

For three processes the file is 31\,MB, and the memory is the same as
without the log (the logged check's time includes growing Sylvan's
tables). Records by operation: image 270K, union 164K, minus 54K; no copy
nodes (a Sylvan feature the plugin does not use), and only the position
kinds copy, read, write and only-read. Those are the cases that need
lemmas.

\subsection{From records to facts}\label{sec:prooflog:facts}

\paragraph{What a fact is.} A \emph{fact} is one lemma applied to its
arguments, whose type the kernel checks. The lemmas about union and minus
are in \file{theories/DdTree.v}:
\begin{itemize}
\item \code{U a b c} (\loc{theories/DdTree.v}{131}): \code{c} holds exactly
  the vectors of \code{a} and those of \code{b};
\item \code{Mi a b R} (\loc{theories/DdTree.v}{183}): \code{R} is within
  \code{a}, and every vector of \code{a} is in \code{b} or in \code{R};
\end{itemize}
and there is one lemma for each case of Sylvan's operations:
\code{u_a0}, \code{u_b0}, \code{u_same} for the cases with an empty or an
equal argument, \code{u_lt}, \code{u_gt}, \code{u_eq}, \code{u_eq0} for the
three comparisons (the last when the downs' union is empty), and
\code{u_lt0}, \code{u_gt0}, \code{u_eq_d}, \code{u_eq_r} for a last step
folded into one fact (\cref{sec:prooflog:fit}); likewise \code{mi_*} for
minus. None of these lemmas needs to know that
$v_a < v_b$: they are about sets, not about order. So the facts of union
and minus need no comparison of numbers at all.

\paragraph{A worked example.} Take $a = \{(1), (3)\}$ =
\code{Nd 1 D1 (Nd 3 D1 D0)} and $b = \{(2)\}$ = \code{Nd 2 D1 D0}. Sylvan's
union compares 1 with 2: $1 < 2$, so the result is \code{Nd 1 D1 R} where
\code{R} is the union of \code{Nd 3 D1 D0} with $b$. That union compares 3
with 2: $3 > 2$, so it is \code{Nd 2 D1 R'} with \code{R'} the union of
\code{Nd 3 D1 D0} with \code{D0}, which is \code{Nd 3 D1 D0} (an empty
argument). The log holds two records (the two compared cases) and the
nodes; the facts, innermost first, are
\begin{verbatim}
f0 := u_b0 (Nd 3 D1 D0)
      : U (Nd 3 D1 D0) D0 (Nd 3 D1 D0)
f1 := u_gt (Nd 3 D1 D0) 2 D1 D0 (Nd 3 D1 D0) f0
      : U (Nd 3 D1 D0) (Nd 2 D1 D0) (Nd 2 D1 (Nd 3 D1 D0))
f2 := u_lt 1 D1 (Nd 3 D1 D0) (Nd 2 D1 D0) (Nd 2 D1 (Nd 3 D1 D0)) f1
      : U a b (Nd 1 D1 (Nd 2 D1 (Nd 3 D1 D0)))
\end{verbatim}
and the last says that the result, $\{(1), (2), (3)\}$, is the union.
(Sylvan stores a union under its own order of the two arguments; the
facts follow the order the caller asked for, each case having its mirror
image, so no fact is needed to turn one around.)

\paragraph{The translation} (\code{Dd_cert.write_log},
\loc{lib/certify/dd_cert.ml}{281}). For each union or minus record, the
translator finds the calls it made from the nodes' contents (which case
applies), and their facts from their own records; it checks that each
record's result is what its case says it must be (an error otherwise:
nothing is trusted), and writes the fact. It computes no set: every
result comes from the log.

\paragraph{Chunks.} A proof of three-process CADP has about 670,000
facts. Built as one term, the whole proof would sit in memory at once
before the kernel checks it. Instead, \code{MeBi Check Log} writes
the facts to a temporary file in pieces called \emph{chunks}, and reads
them back one chunk at a time (\file{lib/certify/dd_read.ml}); a proof's
facts, one chunk, go from the writer to the reader a line at a time,
with no file; each chunk is declared as its own
constant, which the kernel checks when it is declared. Here a chunk holds
up to 20,000 facts and the diagram nodes they mention, as local
definitions of its own proof; its statement is \code{True}, proved by
putting all its facts in a chain (\code{chk_u}, \code{chk_mi},
\code{chk_nil}, \loc{theories/DdTree.v}{264}), so that the kernel checks
every one. (A first design made each node a global constant: for three
processes, about 550,000 node constants took 1.8\,GB before the first chunk
was read, and the memory cap stopped it.)

\paragraph{The command and its numbers.} \code{MeBi Check Log "<file>".}
(\loc{src/user_tactics.ml}{289}) runs the translation and the reading on a
kept log, for testing; it proves no goal. Measured (10 November 2026):

\begin{center}
\small
\begin{tabular}{@{}lrrrrrr@{}}
\toprule
 & chunks & facts & nodes & written & read, checked & memory \\
\midrule
two-process CADP & 1 & 15,200 & 6,183 & 0.01\,s & 0.4\,s & 0.34\,GB \\
three-process CADP & 8 & 668,460 & 581,960 & 1.8\,s & 26.5\,s & 0.83\,GB \\
\bottomrule
\end{tabular}
\end{center}

For three processes, nodes are made again in each chunk that needs them
(233K distinct); the memory grows with each chunk, because Rocq keeps each
declared chunk until it writes the \file{.vo} file. One chunk instead of
eight (\code{MEBI_CHUNK_FACTS}) makes each node once and takes the same
time, but needs 1.42\,GB.

\subsection{Images, complete and sound}\label{sec:prooflog:images}

\paragraph{What they are.} A system's steps are stored as a
\emph{relation}: a diagram of \emph{tuples}, each saying how a step
changes a vector. For each position of the vector, a \emph{code} says what
the relation does there: \emph{copy} (the value stays, and the tuple holds
nothing for it), \emph{only-read} (the value stays and must equal the
tuple's), or \emph{read-write} (the tuple holds the old value, then the new
one). \emph{Example:} with codes (read-write, copy) and the tuple
$(1, 2)$, the vector $(1, 7)$ steps to $(2, 7)$: position 1 must hold 1 and
becomes 2, position 2 is copied. The \emph{image} of a set $x$ through a
relation is every vector one step after a vector of $x$: the image of
$\{(1,7), (3,7)\}$ is $\{(2,7)\}$ ($(3,7)$ has no step: no tuple reads
3). The \emph{inverse} relation swaps each tuple's old and new values:
$(2, 1)$; the image of $\{(2,7)\}$ through it is $\{(1,7)\}$, the vectors
one step \emph{before}. In Rocq the codes are \code{Copy}, \code{RO},
\code{RW} (\file{theories/DdTree.v}, \code{rstep},
\loc{theories/DdTree.v}{65}); Sylvan writes them as numbers per level of
the relation (0, 3, and 1 then 2), and the translator turns one into the
other.

\paragraph{Complete and sound.} A fact about a computed set can say two
things. \emph{Complete}: everything that belongs there is in the result
(the image of every vector of $x$ is in it: \code{PostC},
\loc{theories/DdTree.v}{281}). \emph{Sound}: everything in the result
belongs there (every vector of it is the image of a vector of $x$:
\code{PostS}, \loc{theories/DdTree.v}{526}). The proof needs the mover's
images complete (every move from $z$ is caught) and the answerer's sound
(whatever is in $\mathit{ok}_l$ really can be answered); \code{WlC} and
\code{WlS} are the same for the written part of a read-write position.

\paragraph{Why completeness needs comparisons.} Like union, Sylvan's image
walks the set's and the relation's right chains in parallel and skips a
node whose value cannot match. \emph{Example:} at a read-write position,
the set's first value is 1 and the relation's is 4. Every tuple reads a
value of at least 4 (values increase along a right chain), so no vector
starting with 1 has a step: Sylvan skips the set's node 1. For a sound
fact, skipping only drops candidates, so nothing needs to be shown
(\code{posts_x}, \code{posts_r}). For a complete fact, it must be shown
that what was dropped could not have had a step: that needs $1 < 4$, and
that the chains are sorted (the facts assume the sets sorted, \code{Srt}).
Each skip lemma of the complete direction takes the comparison Sylvan
made as a premise (\code{post_rw_lt}: $v_x < v_r$).

\paragraph{Comparisons as evidence.} A comparison such as $5 < 12$ is not
left for the kernel to compute: it is proved from the two numbers' binary
digits, one lemma per digit (\code{plt_*}, \code{nlt_pos}). In Rocq a
positive number is a list of binary digits read from the last:
$5 = $ \code{xI (xO xH)} (101), $12 = $ \code{xO (xO (xI xH))} (1100).
The proof of $5 < 12$ compares the last digits first: \code{plt_xI_xO}
says $2p+1 < 2q$ when $p < q$, so it needs $2 < 6$; \code{plt_xO_xO} says
$2p < 2q$ when $p < q$, so it needs $1 < 3$; and \code{plt_1_xI} says
$1 < 2q+1$. Three lemma applications; a number of $n$ binary digits needs
at most $n$.

\paragraph{When the relation ends early.} A relation of one system
leaves the other system's positions alone, and a relation touches only
some of its own. Its codes stop at its last touched position, and
\code{rstep} carries the positions after the last code unchanged, as
Sylvan does past the end of its meta. So an image at that point is the set
itself, in one fact (\code{post_end}, \code{posts_end},
\loc{theories/DdTree.v}{298}), with no step per node below it.

\paragraph{Related work.} Producing proof steps from a decision-diagram
package as it computes is the idea of Sinz and Biere~\cite{sinz-biere06},
and of Bryant and Heule~\cite{bryant-heule21}, where each BDD operation
emits clauses of an extended-resolution proof. Here each LDD operation's
case becomes the application of a lemma about sets, checked by Rocq's
kernel.

\subsection{Making three-process CADP fit}\label{sec:prooflog:fit}

Three-process CADP (7.9 million states) is the largest example. Each
change below was measured with \file{test/cert_tool.exe --log}, which runs
the translator alone and counts the facts of each lemma.

\paragraph{1. The check as images.} The game's own check, logged, needed
the explored pairs as a universe, each system's exploration closed,
intersections with the universe and complete preimages: 1.38M records,
2.46M facts. The forward check (\cref{sec:prooflog:background}) needs
none of these: 488K records.

\paragraph{2. The answering side first.} A pair is one vector. A relation
of the side whose positions come second must copy the other side's
positions before it reaches its own, one fact per node (\emph{copy
levels}). The answerer's relations are applied most (each label, each
silent round), so its positions come first (\loc{lib/learn/learn.ml}{1305};
for a bisimilarity, where both answer, the narrower side's), and each
meta stops at its relation's last touched position (\code{meta_levels},
\loc{lib/learn/learn.ml}{520}). The answerer's images are about 100 facts
in all.

\paragraph{3. Each image one way only: its role.} The logged run records
each image it makes at its top level with the way the proof uses it, 1
sound or 2 complete (\code{Ldd.log_top_relprod} in \file{lib/ldd}, the
fork's \code{lddmc_log_top}, operation 6 in the log), with the arguments
as passed, before Sylvan skips any node. The translator proves each image
only in that direction; a sound image needs no comparison.

\paragraph{4. Smaller tables in the translator.} The translator looks up a
record by its call (the operation and its arguments) in a table of its own
(\code{Itab} in \file{lib/certify/dd_cert.ml}: open addressing, the call
packed into two integers, the result and a ``called by another record''
bit in a third), about 32 bytes a record, and reads the log as a stream
(\code{Proof_log.read_iter}).

\paragraph{5. Comparisons shared within a chunk.} A complete skip needs the
comparison \code{a < b} of two node values, proved digit by digit
(\cref{sec:prooflog:images}). The same pair of values recurs: its proof is
made once per chunk (three processes: about 17,000 comparison facts).

\paragraph{6. A chain's last step folded in.} Many facts end a chain with
a fact about an empty set: the last written value of a read-write
position, then the empty rest (\code{wl_nd}, then \code{wl_r0}); a node
with no right sibling (\code{post_copy}, then \code{post_x0}); a union
with an empty or an equal side. A lemma for each such pair, proved once
from the two (\code{wl_last}, \code{post_copy_last}, \code{post_rw_last},
\code{u_lt0}, \code{u_gt0}, \code{u_eq_d}, \code{u_eq_r}, \code{mi_eq0d},
\code{mi_eq0r}), makes it one fact; and each union is proved in the order
of its arguments the caller asked for, with no fact to swap them. Three
processes: 963,941 facts down to 668,460, and 10\% less time (the folded
facts were the small ones).

\paragraph{7. The largest chunk first.} Rocq's kernel, when it checks a
constant, first copies its whole term into a hash-consed form (each
distinct subterm once, with a table), and keeps the types of the shared
parts. Measured on a chunk of 292,000 facts: its term took 52\,MB, and
checking it 510\,MB more, most of it live data, not garbage (a GC four
times as eager still needed 386\,MB). So a chunk's check needs memory in
proportion to its size, on top of what the chunks declared before it keep
(about 400 bytes a fact). One root alone can give 300,000 facts (one
image of the game's check), so chunks cannot simply be made smaller. Since the
check chunks need nothing from each other, \code{Dd_read} reads them
largest first: the largest check then meets the least memory held.
(Cutting a large root across chunks would need the nodes shared between
chunks as constants of their own; measured, a constant costs about 3\,KB,
too much for 233,068 nodes.)

\paragraph{Where the time goes.} Measured on three processes: the kernel
spends about 14\,\textmu s per node \code{let} and 28\,\textmu s per fact,
typing and hash-consing about equally. Binding the lists of codes and the
numbers as \code{let}s, or one chunk instead of eight, does not change the
total: what reduces it is fewer facts and fewer nodes.

\paragraph{Measured} (10 November 2026, \file{bench/prooflog.py}):
\begin{center}
\small
\begin{tabular}{@{}lrrrr@{}}
\toprule
 & decided & facts & checked in & memory at most \\
\midrule
two-process CADP & 0.8\,s & 15,200 & 0.7\,s & 0.34\,GB \\
three-process CADP & 22.9\,s & 668,460 & 29.5\,s & 0.83\,GB \\
\bottomrule
\end{tabular}
\end{center}
Before these changes, three processes took 2,461,163 facts, checked in
121\,s, 1.69\,GB at most. These are the facts alone (\code{MeBi Check
Log}); \cref{sec:qed} builds the proof of the goal from them.

\subsection{Trying it}\label{sec:prooflog:try}

For the two CADP examples, \file{bench/prooflog.py} does all of this
(after \code{dune build}): it copies each file into \file{_certify/log/},
turns the log on, stops once the engine has decided, checks each log in a
file of its own next to it, and prints the facts, times and memory. A
complete example is \file{test/rocq/LogTest.v} (in the build). In short:
\begin{verbatim}
Require Import MEBI.loader.
(* ... an LTS [step], then: *)
Set MeBi Log Mode "keep".
Example logged : weak_bisimilar step step (par A B) (par B A).
Proof. weak_bisim. Qed.
Set MeBi Log Mode "off".
MeBi Check Log ".LogTest.v.mebi/logged.log".
Print Assumptions mebi_log1_c1.
\end{verbatim}
\begin{enumerate}
\item The tactic decides the goal and proves it from the certificate
  (\cref{sec:qed}); with \code{"keep"}, the engine also leaves
  \file{.LogTest.v.mebi/logged.log}, and prints on standard error a line
  \code{[proof log] 628 records; the check as images: passes \dots}: the
  number of records, and that the check holds in the repeated run (it
  must pass).
\item \code{MeBi Check Log} prints the number of chunks, facts and nodes,
  the times, and the name of the last chunk's constant
  (\code{mebi_log1_c1}: the first use of the command in the file, its
  first chunk).
\item \code{Print Assumptions} on that constant must say ``Closed under the
  global context'': the facts use no axiom.
\end{enumerate}
To measure a large example, add \code{Set MeBi Log Mode "keep".} before its
proof and run it with \code{MEBI_CERT_STOP=decided} (stop once the engine
has decided) and \code{MEBI_CERT_STATS=1}; then check the log in a second
file. Run heavy examples under a memory cap (\file{CLAUDE.md}).

\subsection{What remains, and what was decided}\label{sec:prooflog:next}

\paragraph{Next} (\file{doc/certificate.md}; \cref{sec:status}): the
optimisations (\cref{sec:optimisations}); every example's proof already
ends with \code{Qed} (\cref{sec:qed:cost}).

\paragraph{Decided} (10 November 2026):
\begin{itemize}
\item \emph{one extra logged check}, not a log of every iteration;
\item \emph{the logged check as images}, the mover's complete and the
  answerer's sound by inverse relations, not the game's backward check:
  no universe, no closure of the exploration, no intersection, no
  complete preimage;
\item \emph{evidence, not re-checking}: where Sylvan has compared two
  values, the comparison is written as a proof built from the two
  numbers' binary digits, one lemma per digit, not as a computation the
  kernel repeats; lookups in lists are avoided by giving the element
  itself;
\item \emph{steps from the rules only}: the engine finds a term's steps by
  interpreting the LTS's rules; a step it cannot follow is an error;
\item \emph{side conditions with a \code{Decidable} instance}: proved
  through the instance, the kernel reducing its boolean once per step.
\end{itemize}
